\chapter{La sexualidad como don}

\epigraph{El cuerpo, y solo él, es capaz de hacer visible lo que es invisible: lo espiritual y lo divino.}{Juan Pablo II}

Hay pocos temas sobre los que se hable tanto y se diga tan poco.

La sexualidad humana está en todas partes ---en la publicidad, en las series, en las conversaciones de bar, en los titulares--- y, sin embargo, rara vez se habla de ella con la profundidad que merece. Nos movemos entre dos extremos que parecen opuestos pero que comparten una misma raíz: no entender qué es. De un lado, la represión: la sexualidad como algo sucio, peligroso, de lo que es mejor no hablar. Del otro, la banalización: la sexualidad como consumo, como entretenimiento, como un acto más entre otros que no compromete a nada ni a nadie.

Entre el tabú y el escaparate hay un abismo. Y en ese abismo se ha perdido algo esencial. Algo que nuestros cuerpos saben pero que la cultura ha olvidado. Algo que la fe cristiana custodia, no como prohibición, sino como buena noticia.

Vamos a intentar recuperarlo.

\section{Ni represión ni consumo}

Empecemos por desmontar los dos extremos, porque ambos hacen daño.

La represión tiene una historia larga. Durante siglos, una cierta lectura del cristianismo ---más influida por el neoplatonismo y el jansenismo que por el Evangelio--- presentó la sexualidad como una concesión a la debilidad de la carne, algo tolerado en el matrimonio solo porque servía para tener hijos, pero sospechoso en sí mismo. El placer era visto con desconfianza. El cuerpo, como cárcel del alma. El deseo, como enemigo.\footnote{Esta visión negativa del cuerpo y la sexualidad debe más al dualismo platónico y a corrientes rigoristas que al mensaje bíblico. El Cantar de los Cantares, incluido en el canon de las Escrituras, es un canto apasionado al amor erótico entre los esposos. La tradición de la Iglesia siempre ha afirmado la bondad de la creación ---incluido el cuerpo y la sexualidad---, aunque no siempre su predicación haya estado a la altura de esa afirmación.}

El resultado fue generaciones enteras para las que hablar de sexualidad equivalía a hablar de pecado. Y muchos que crecieron así ---quizá vuestros padres, quizá vosotros mismos--- arrastran heridas reales: vergüenza, culpa, incapacidad de vivir la intimidad con libertad y alegría.

En el otro extremo, la cultura contemporánea ha hecho algo distinto pero igualmente empobrecedor: ha vaciado la sexualidad de significado. Si antes el cuerpo era sospechoso, ahora es irrelevante. El sexo se presenta como una actividad recreativa, un intercambio de sensaciones entre cuerpos que se usan mutuamente, sin que el acto diga nada ni comprometa a nada. «Es solo sexo», se dice, como si esas tres palabras no fueran una contradicción en sí mismas.

Porque el sexo nunca es «solo» sexo. Cualquiera que haya vivido la experiencia lo sabe. La intimidad sexual nos deja vulnerables de una forma que ninguna otra actividad humana consigue. Nos expone. Nos desnuda ---y no hablo solo de la ropa---. Hay algo en el acto sexual que toca una fibra profunda, algo que no se puede reducir a biología ni a deporte.

Los dos extremos, curiosamente, comparten el mismo error: tratar la sexualidad como si fuera \emph{solo} un asunto del cuerpo. Para la represión, el cuerpo es el problema. Para la banalización, el cuerpo es el instrumento. Ninguno de los dos entiende que el cuerpo es el \emph{lenguaje}.

\section{Lo que el cuerpo dice}

En los capítulos anteriores hemos visto que el cuerpo humano no es un accidente biológico ni una máquina indiferente. El cuerpo habla. Tiene su propia gramática, su propia elocuencia. La oxitocina que se dispara con el contacto piel con piel, las hormonas que empujan al vínculo, el rostro que se sonroja, las manos que tiemblan cuando rozan por primera vez las de la persona amada: todo eso es lenguaje. El cuerpo dice cosas que la boca a veces no sabe decir.

Pues bien: si hay un momento en que el cuerpo habla con toda su elocuencia, ese momento es la unión sexual.

Juan Pablo II dedicó años a pensar sobre esto. Entre 1979 y 1984, en una serie de audiencias generales que pasaron relativamente desapercibidas, desarrolló lo que se conoce como la \emph{Teología del Cuerpo}: una reflexión sobre el significado del cuerpo humano y la sexualidad que muchos consideran una de las aportaciones más originales de la teología del siglo XX.\footnote{Juan Pablo II, \textit{Hombre y mujer los creó. Catequesis sobre el amor humano}, 129 audiencias generales entre septiembre de 1979 y noviembre de 1984. La obra completa ha sido editada en múltiples volúmenes. Para una introducción accesible en español, cf. Yves Semen, \textit{La sexualidad según Juan Pablo II}, Desclée de Brouwer, 2005.} No es un texto fácil ---su estilo es denso, filosófico, a ratos casi impenetrable---, pero la idea central es de una sencillez luminosa.

La idea es esta: el cuerpo humano tiene un «significado nupcial». Es decir, está hecho para la entrega. No solo puede entregarse: está \emph{diseñado} para ello. Toda la configuración del cuerpo humano ---la diferencia sexual, la complementariedad, el deseo que nos empuja hacia el otro--- apunta a un mismo destino: la donación de uno mismo a la persona amada.

Y esa donación alcanza su expresión más plena, más radical, más elocuente, en la unión sexual de los esposos.

¿Qué dice el cuerpo cuando dos personas se unen en el acto sexual? Dice algo enorme. Dice algo que, si nos paramos a escucharlo, resulta casi sobrecogedor.

El cuerpo dice: \emph{Te doy todo lo que soy.}

No una parte. No un rato. No con condiciones. \emph{Todo.} La desnudez física es signo de una desnudez más profunda: me muestro ante ti tal como soy, sin máscaras, sin defensas, sin reservas. Me entrego entero. Te acojo entera.

Por eso la sexualidad no puede ser trivial. Por eso «solo sexo» es una frase imposible. Porque el cuerpo no sabe mentir ---o, mejor dicho, cuando miente, algo se rompe por dentro---. Cuando dos cuerpos se unen diciendo «te doy todo» mientras la voluntad piensa «te doy un rato», hay una fractura. Una incoherencia entre lo que el cuerpo proclama y lo que la persona decide. Y esa fractura, aunque no siempre sea consciente, deja huella.

\section{Las tres palabras del cuerpo}

Aquí es donde todo lo que hemos visto en los capítulos anteriores se hace carne. Literalmente.

Hemos hablado de las tres propiedades del matrimonio cristiano: indisolubilidad, fidelidad, fecundidad.\footnote{Cf. capítulos 8, 9 y 10 de este libro.} Las hemos visto como regalos, como expresiones del amor total. Pues bien: en la unión sexual, esas tres propiedades dejan de ser conceptos teológicos y se convierten en lenguaje del cuerpo. Lo que los esposos se prometieron con palabras el día de la boda, sus cuerpos lo repiten cada vez que se unen.

\subsection*{El cuerpo dice: «Para siempre»}

Cuando dos esposos se unen, sus cuerpos dicen indisolubilidad. La entrega sexual es, por su propia naturaleza, una entrega sin fecha de caducidad. No es un préstamo. No es un alquiler. Es un don total, sin reservas, sin letra pequeña que diga «hasta que me canse» o «mientras me convenga».

«Serán una sola carne»,\footnote{Gn 2,24.} dice el Génesis. No «una sola carne de momento». No «una sola carne a tiempo parcial». Una sola carne, sin más. La unión de los cuerpos dice para siempre con una elocuencia que las palabras no alcanzan. Cada vez que los esposos se unen, renuevan con sus cuerpos la promesa que hicieron con sus labios: \emph{No me voy. Me quedo contigo. Pase lo que pase.}

\subsection*{El cuerpo dice: «Solo a ti»}

La entrega sexual dice también fidelidad. Y no la fidelidad entendida como restricción ---«no puedo estar con nadie más»--- sino como concentración de toda la capacidad de donación en una sola persona: «todo lo que soy, para ti; todo lo que tengo, para ti; toda mi capacidad de amar, volcada en ti».

Es exclusiva no por prohibición, sino por lógica interna. Si me doy \emph{entero} a ti, no me queda nada para dar a otro. No porque alguien me lo prohíba, sino porque así funciona la entrega total: es indivisible. No se puede dar todo a dos personas a la vez, igual que no se puede estar entero en dos sitios al mismo tiempo.

Los cuerpos que se unen dicen \textit{fides}, fe: «Creo en ti. Me fío de ti. Deposito en tus manos lo más íntimo que tengo ---mi cuerpo, mi vulnerabilidad, mi desnudez--- porque confío en que lo vas a custodiar.»\footnote{Cf. capítulo 9 de este libro, donde se desarrolla la fidelidad como fe depositada en el otro.}

\subsection*{El cuerpo dice: «Estamos abiertos a lo que venga»}

Y los cuerpos que se unen dicen fecundidad. Pero aquí conviene hacer una pausa, porque en el capítulo anterior hemos visto que la fecundidad del matrimonio es algo mucho más amplio que tener hijos. Un matrimonio fecundo es el que irradia vida a su alrededor, el que es sal de la tierra, el que genera esa envidia sana que hace que otros deseen amar así. Los hijos, cuando los hay, son parte de esa fecundidad ---una parte hermosa y central---, pero no la agotan.

Lo que la unión sexual expresa es una dimensión específica de esa fecundidad amplia: la apertura a la vida biológica. Los esposos no se unen «a pesar de» la posibilidad de engendrar, ni «para» engendrar como si fuera una producción industrial. Se unen como don total ---y un don total incluye la fertilidad, incluye la posibilidad de que ese amor sea tan real, tan concreto, que tome carne en una persona nueva---.

La Teología del Cuerpo de Juan Pablo II desarrolló esta idea con una profundidad que no tenía precedentes.\footnote{Cf. Juan Pablo II, \textit{Hombre y mujer los creó}, audiencias 117-133; Pablo VI, Encíclica \textit{Humanae Vitae} (1968), especialmente nn. 11-12, donde se desarrolla la inseparabilidad del significado unitivo y procreativo del acto conyugal. También Juan Pablo II, \textit{Familiaris Consortio} (1981), nn. 28-35.} Para él, si la entrega sexual es total, si los cuerpos dicen «te doy todo lo que soy», entonces cerrar deliberadamente la puerta a la procreación es introducir una reserva en el lenguaje del cuerpo. Es decir «te doy todo» mientras se retiene algo.

Conviene no confundir dos planos distintos. Como vimos en el capítulo anterior, la fecundidad global de un matrimonio es inmensamente amplia: abarca lo social, lo espiritual, el testimonio, la acogida. Un matrimonio sin hijos puede ser extraordinariamente fecundo. Pero aquí no hablamos de la fecundidad de la vida entera, sino del significado de un gesto concreto: el acto de unión conyugal. Y ese gesto tiene su propia gramática. No se trata de que cada encuentro deba producir un hijo ---eso nadie lo sostiene---, sino de que el lenguaje que hablan los cuerpos al unirse sea el de una entrega sin letra pequeña, donde la apertura a la vida es una dimensión de la donación total, no una obligación de resultado.

Es una enseñanza exigente, y conviene recibirla con honestidad: es la visión de san Juan Pablo II, coherente con toda su antropología del cuerpo. Pero probablemente la fecundidad a la que está llamado el matrimonio se pueda ---y se deba--- ver como un todo más amplio que la procreación. La procreación no queda excluida de esa visión más amplia: queda \emph{incluida} en ella, como una de sus expresiones más bellas, pero no la única. Un matrimonio sin hijos no es un matrimonio estéril si su amor genera vida a su alrededor.\footnote{Cf. capítulo 10 de este libro, donde se desarrolla la fecundidad como irradiación. También la exhortación del Papa Francisco: «El amor siempre da vida. [...] El amor conyugal no se agota dentro de la pareja [...] sino que los habilita para la máxima donación.» Francisco, \textit{Amoris Laetitia}, n. 165 (2016).} Y un matrimonio con hijos, si vive replegado sobre sí mismo, puede estar traicionando una fecundidad que debería desbordarlo.

La apertura a la vida ---en su sentido más pleno--- no es una carga impuesta desde fuera. Es la consecuencia natural de un amor que se da sin reservas. Es lo que pasa cuando el «te doy todo» se dice de verdad, con el cuerpo y con la vida entera. Y esa vida que nace del amor toma muchas formas: a veces tiene la cara de un hijo, a veces la de un hogar acogedor, a veces la de un servicio callado al que lo necesita.

\section{La castidad como integración}

Si hay una palabra que la cultura contemporánea ha malentendido más que ninguna otra, esa palabra es \emph{castidad}.

Para la mayoría de la gente, castidad significa abstinencia. No tener relaciones sexuales. Punto. Y como la abstinencia no suele ser un proyecto de vida demasiado atractivo, la castidad ha quedado relegada al rincón de las virtudes polvorientas, de las que nadie habla excepto para burlarse.

Pero castidad no es abstinencia. La castidad es algo mucho más amplio, más dinámico, más vivo. La castidad es \emph{integración}.\footnote{«La castidad significa la integración lograda de la sexualidad en la persona, y por ello en la unidad interior del hombre en su ser corporal y espiritual.» \textit{Catecismo de la Iglesia Católica}, n. 2337.}

¿Qué significa eso? Significa que la sexualidad no es una fuerza salvaje que hay que encerrar en una jaula para que no haga daño, ni una energía neutra que se puede gastar como uno quiera sin consecuencias. La sexualidad es una dimensión profunda de la persona ---tan profunda como la inteligencia o la libertad--- que necesita ser integrada en el conjunto de lo que somos. Ordenada. No reprimida: \emph{ordenada}. Puesta al servicio del amor, que es para lo que existe.

Una persona casta no es una persona que no desea. Es una persona que desea \emph{bien}. Que ha aprendido a mirar al otro como persona y no como objeto. Que ha integrado su deseo sexual en su capacidad de amar, de modo que la atracción no le ciega sino que le ilumina, no le esclaviza sino que le libera para amar mejor.

La castidad es, para los esposos, la capacidad de vivir su sexualidad como lenguaje de donación y no de apropiación. De unirse al otro para \emph{darse}, no para \emph{tomarse}. De buscar en la intimidad el bien del otro antes que la propia satisfacción ---lo cual, paradójicamente, es el camino más seguro hacia una intimidad más plena y más gozosa---.

No es fácil. Nadie dijo que lo fuera. Pero es profundamente liberador, porque transforma el deseo de fuerza centrípeta ---todo hacia mí--- en fuerza centrífuga ---todo hacia ti---. Y eso es exactamente lo que hace el amor cuando es amor de verdad.

\section{La escuela de todos los días}

Y aquí quiero hablar de algo que va más allá de la sexualidad, pero que nace de ella. Porque la entrega que los esposos aprenden en la intimidad no se queda entre las sábanas. Se extiende, como una marea, a toda la vida compartida.

El matrimonio transforma. No deja a nadie igual.

Pensadlo un momento: no sois la misma persona que erais cuando os casasteis ---o cuando decidisteis casaros---. Y no me refiero solo a las canas o a los kilos. Me refiero a algo más profundo. La convivencia diaria con otra persona, la entrega concreta y cotidiana, la fricción inevitable de dos vidas que se funden sin dejar de ser distintas\ldots{} todo eso os ha ido cambiando. Os ha ido puliendo. Os ha ido haciendo ---si os habéis dejado--- mejores personas de las que habríais sido solos.

Porque el matrimonio es una escuela. No una escuela con pupitres y exámenes, sino una escuela de \emph{virtudes}, de esas que se aprenden en la piel, no en los libros.

Se aprende \textbf{humildad}, porque tarde o temprano ---más temprano que tarde--- hay que reconocer que te has equivocado. Que has dicho algo injusto. Que has sido egoísta. Y pedir perdón. Y el perdón pedido a la persona que duerme a tu lado es más costoso y más verdadero que cualquier examen de conciencia abstracto.

Se aprende \textbf{paciencia}, porque el otro no es como tú querrías que fuera. Tiene sus manías, sus ritmos, sus silencios que no entiendes y sus palabras que sobran. Y aprender a convivir con todo eso sin pretender cambiar al otro a tu imagen y semejanza es un ejercicio de paciencia que ningún libro puede enseñar.

Se aprende \textbf{justicia}, porque hay que ceder. Hay que repartir. Hay que escuchar la perspectiva del otro y reconocer que quizá ---solo quizá--- tiene razón. La vida compartida exige un sentido de la justicia que no se aprende en la facultad de Derecho sino en la cocina, en la logística de quién lleva a los niños al colegio, en la negociación silenciosa de cada día.

Se aprende \textbf{generosidad}, porque amar es dar sin llevar cuentas. Dar tiempo, dar atención, dar presencia, dar renuncia. Dar sin esperar retorno, sin pasar factura, sin el cálculo mezquino de «yo hice esto, ahora te toca a ti».

Y se aprende \textbf{fortaleza}, porque hay momentos en que todo es difícil. En que el cansancio pesa más que el cariño. En que la rutina ha borrado el brillo. En que uno querría salir corriendo. Y decidir quedarse ---no por inercia, no por cobardía, sino por amor--- es un acto de fortaleza que deja huella en el alma.

No es que el matrimonio enseñe estas virtudes en abstracto, como quien lee un catálogo de buenas intenciones. Es que la convivencia diaria, el roce constante con otra persona, te las va tallando en la carne. Te pule, te lima, te moldea. Te hace mejor persona ---si te dejas---. Y ese «si te dejas» es importante, porque la escuela del matrimonio solo funciona cuando uno acepta ser alumno. Cuando uno reconoce que no ha llegado terminado, que le queda mucho por aprender, que el otro es maestro tanto como compañero.

El matrimonio no solo une a dos personas. Las transforma. Y la persona que sale de esa transformación no es la que entró, ni tampoco la que habría sido si hubiera caminado sola.

\section{El sacramento hecho carne}

Volvamos a la sexualidad para cerrar el círculo.

Decíamos al principio que el cuerpo habla, y que en la unión sexual habla con su máxima elocuencia. Ahora podemos ver el cuadro completo: la sexualidad de los esposos es el lugar donde el sacramento del matrimonio se hace carne de la forma más literal posible.

Cuando los esposos se casaron, dijeron tres «sí» ante el altar: sí a la indisolubilidad, sí a la fidelidad, sí a la fecundidad.\footnote{El rito del matrimonio incluye tres preguntas a los contrayentes antes del consentimiento: si vienen libremente, si se comprometen a amarse y respetarse durante toda la vida, y si están dispuestos a recibir con amor los hijos que Dios les dé. Cf. \textit{Ritual del Matrimonio}, Praenotanda, nn. 2-3.} Esos tres «sí» no fueron solo palabras. Fueron el compromiso de una vida entera.

Y cada vez que se unen, sus cuerpos repiten esos tres «sí»: me doy a ti para siempre, me doy a ti en exclusiva, me doy a ti abierto a la vida que pueda venir. La unión sexual de los esposos no es un apéndice del sacramento, un añadido más o menos relevante. Es el sacramento \emph{en acto}. Es la alianza renovada en la carne. Es el «sí» del altar traducido al lenguaje del cuerpo.

Por eso la sexualidad conyugal es sagrada. No sagrada en el sentido de intocable o de solemne. Sagrada en el sentido más profundo: es un lugar de encuentro con el misterio. Un lugar donde dos personas, en su vulnerabilidad y su ternura, hacen visible algo del amor de Dios. Ese Dios que se entrega sin reservas, que es fiel hasta el final, y cuyo amor es tan fecundo que crea vida.

No hacen falta velas ni música de fondo para que esto sea verdad. Es verdad en la intimidad de un martes por la noche, con el cansancio de la jornada encima y los niños por fin dormidos. Es verdad en la ternura torpe de las primeras veces y en la complicidad callada de los muchos años. Es verdad siempre que dos esposos se unen con la intención de darse, no de tomarse; de amar, no de usar.

Y esa verdad ---que el cuerpo habla, que la sexualidad es don, que la intimidad conyugal dice cosas que las palabras no alcanzan--- es quizá la mejor noticia que la fe cristiana tiene que ofrecer en un mundo que ha olvidado escuchar lo que el cuerpo dice.

\paracaminar

\begin{enumerate}
  \item ¿Os reconocéis en alguno de los dos extremos ---represión o banalización--- a la hora de entender la sexualidad? ¿Cómo ha influido vuestra educación, vuestra familia o vuestra cultura en la forma en que vivís la intimidad?

  \item Pensad en las virtudes que el matrimonio os ha ido enseñando ---o que intuís que os enseñará---. ¿Cuál os cuesta más? ¿En cuál habéis crecido más?

  \item La idea de que el cuerpo «habla» en la unión sexual, y que dice indisolubilidad, fidelidad y apertura a la vida: ¿cómo la recibís? ¿Cambia en algo vuestra forma de entender la intimidad conyugal?
\end{enumerate}

\vinetacapitulo{ch11}
